% ---------------------------------------------------------------
% ICML 2026 format — Project Update 2: Data & Formal Definitions
%
% HOW TO COMPILE IN OVERLEAF:
%   1. Open the ICML 2025 Overleaf template and copy icml2025.sty:
%      https://www.overleaf.com/latex/templates/icml2025-template/dhxrkcgkvnkt
%      Rename it icml2026.sty and upload to your Overleaf project.
%   2. Upload icml2026.sty + this .tex + your references.bib
%   3. Compile with pdflatex
% ---------------------------------------------------------------
\documentclass{article}
\usepackage[accepted]{icml2025} 
\usepackage[numbers]{natbib}

\usepackage{amsmath}
\usepackage{booktabs}
\usepackage{graphicx}
\usepackage{url}
\usepackage[hidelinks]{hyperref}
\usepackage{tikz}
\usepackage{float}
\usetikzlibrary{shapes.geometric, arrows.meta, positioning, backgrounds, fit}

% ---------------------------------------------------------------
\icmltitlerunning{Corporate Credit Card Misuse Detection}

\begin{document}

\twocolumn[
\icmltitle{Corporate Credit Card Misuse Detection Using\\
           Machine Learning and ETL Pipeline\\[4pt]
           {\normalsize\rm Project Update 2 --- Data \& Formal Definitions}}

\begin{icmlauthorlist}
\icmlauthor{Pooja Malakappa Nuchchi}{neu}
\end{icmlauthorlist}

\icmlaffiliation{neu}{Khoury College of Computer Sciences, Northeastern University (Roux Institute), Portland, ME, USA}

\icmlcorrespondingauthor{Pooja Malakappa Nuchchi}{nuchchi.p@northeastern.edu}

\icmlkeywords{fraud detection, machine learning, ETL pipeline, class imbalance, SMOTE}

\vskip 0.3in
]

\printAffiliationsAndNotice{}

% ---------------------------------------------------------------
\begin{abstract}
Corporate credit card misuse is a significant financial and compliance
challenge in enterprise settings. This paper presents an end-to-end
machine learning system for detecting fraudulent corporate card
transactions, combining a full ETL pipeline (PostgreSQL, Apache Airflow)
with a five-model comparison (Logistic Regression, LinearSVC, Random
Forest, XGBoost, LightGBM). We use the ULB Credit Card Fraud Detection
dataset (284,807 transactions, 0.17\% fraud) and address severe class
imbalance using SMOTE. Results are served through a FastAPI backend and
React dashboard for production-style deployment.
\end{abstract}

% ---------------------------------------------------------------
\section{Introduction and Problem Statement}

Corporate credit card misuse is a significant financial and compliance
challenge. Employees issued corporate cards for legitimate business
expenses may use them for unauthorized personal purchases, resulting in
financial losses and policy violations. Manual review of thousands of
transactions is time-consuming and not scalable.

Formally, the learning problem is defined as: given a transaction record
$x \in \mathbb{R}^{30}$ consisting of 28 anonymized behavioral features
$V_1, \ldots, V_{28}$, a normalized transaction amount, and hour of day,
predict a binary label $y \in \{0, 1\}$ where $y = 1$ indicates misuse
and $y = 0$ indicates legitimate spending. The dataset used is the Credit
Card Fraud Detection dataset (ULB, Kaggle) containing 284,807 transactions
with 99.83\% legitimate vs.\ 0.17\% fraudulent --- a severely imbalanced
classification problem.

% ---------------------------------------------------------------
\section{Motivation}

Financial fraud detection is one of the most practically important
applications of machine learning. Organizations lose significant revenue
to fraud annually, with corporate card misuse among the most common types.
Automated detection systems reduce investigation time and improve accuracy
compared to rule-based systems. Multiple studies have shown that models
like Random Forest and Logistic Regression can effectively detect
fraudulent transactions~\cite{wang2023credit,hande2019credit}. The core
challenge is severe class imbalance, which requires techniques like SMOTE
oversampling and data augmentation~\cite{khalid2024enhancing,strelcenia2023improving}.

% ---------------------------------------------------------------
\section{Related Work}

\textbf{Logistic Regression} is a common baseline for fraud
classification because it is simple and interpretable. Wang et al.\
compared Logistic Regression and Random Forest on simulated transaction
data and found Random Forest performs better~\cite{wang2023credit}.
Sopiyan et al.\ also compared Logistic Regression with Random Forest and
Gradient Boosting using SMOTE and found ensemble methods outperform the
logistic baseline~\cite{sopiyan2022fraud}.

\textbf{Support Vector Machines} have been used in fraud detection for
finding decision boundaries in high-dimensional spaces. Xia et al.\
applied SVM with GridSearchCV tuning to the Kaggle credit card dataset,
optimizing AUC and F1-score~\cite{xia2022credit}. Aftab et al.\ used SVM
alongside Random Forest and Logistic Regression and found ensemble methods
outperform single classifiers~\cite{aftab2023fraud}.

\textbf{Random Forest} is one of the most effective models for fraud
detection because its ensemble voting mechanism naturally handles class
imbalance. Hande et al.\ compared Logistic Regression, Decision Tree, and
Random Forest and found Random Forest performs best~\cite{hande2019credit}.
Aburbeian and Ashqar combined Random Forest with SMOTE and achieved strong
results on imbalanced credit card data~\cite{aburbeian2023credit}.

\textbf{Gradient Boosting} methods are the current state of the art for
tabular fraud detection. XGBoost builds trees sequentially, each
correcting the residual errors of the previous. LightGBM uses leaf-wise
tree growth, which makes it faster on large datasets. Sopiyan et al.\
found Gradient Boosting competitive with Random Forest on credit card
fraud data~\cite{sopiyan2022fraud}.

\textbf{Class Imbalance Handling:} SMOTE creates synthetic minority class
samples by interpolating between existing ones. Khalid et al.\ tested
ensemble methods with SMOTE on European credit card transactions and showed
that resampling improves minority class detection~\cite{khalid2024enhancing}.
Tran and Dang compared SMOTE and ADASYN for balancing credit card fraud
data before training Random Forest and Logistic
Regression~\cite{tran2021machine}. Strelcenia and Prakoonwit proposed
K-CGAN as an alternative to SMOTE~\cite{strelcenia2023improving}.

\textbf{Production ML Pipelines:} Most prior studies focus only on model
accuracy without considering deployment. ETL pipelines that separate
ingestion, transformation, and loading ensure reproducibility and
scalability. This project adopts that principle using PostgreSQL and
Apache Airflow.

% ---------------------------------------------------------------
\section{How This Approach Differs}

Most prior work focuses only on model performance in isolation. This
project combines a complete data engineering layer (ETL pipeline,
PostgreSQL, Apache Airflow) with a five-model comparison (Logistic
Regression, LinearSVC, Random Forest, XGBoost, LightGBM), a REST API
backend (FastAPI), and an interactive React dashboard. This end-to-end
production-style system addresses real-world deployment considerations
that are typically missing from academic fraud detection papers.


\begin{figure}[H]
\centering
\includegraphics[width=\columnwidth]{architecture.png}
\caption{End-to-end system architecture. Raw transaction data flows
through the ETL pipeline orchestrated by Airflow, is balanced with
SMOTE, then passed to five ML models. Random Forest ($\star$,
AUC-ROC~0.9845) is served via FastAPI and visualized in a React
dashboard.}
\label{fig:arch}
\end{figure}

\caption{End-to-end system architecture. Raw transaction data flows
  through the ETL pipeline orchestrated by Airflow, is balanced with
  SMOTE, then passed to five ML models. Random Forest ($\star$,
  AUC-ROC~0.9845) is served via FastAPI and visualized in a React
  dashboard.}
\label{fig:arch}
\end{figure}

% ---------------------------------------------------------------
\section{Dataset}

\subsection{Input and Output}

The input is a transaction record $x \in \mathbb{R}^{30}$ with three
types of features: (1) 28 PCA-transformed anonymized features
$V_1, \ldots, V_{28}$ representing behavioral patterns,
(2)~\texttt{Amount\_Scaled}, the transaction amount standardized using
z-score normalization:
\[
  z = \frac{x - \mu}{\sigma}
\]
and (3)~\texttt{Hour}, the hour of day extracted from the raw
\texttt{Time} field as:
\[
  \text{Hour} = \left\lfloor \frac{\texttt{Time}}{3600} \bmod 24
  \right\rfloor
\]
The output is a binary label $y \in \{0, 1\}$ where $y = 1$ indicates
fraud.

\subsection{Distribution, Size, and Missingness}

The dataset contains 284,807 transactions. The class distribution is
severely imbalanced: 284,315 legitimate (99.83\%) and 492 fraudulent
(0.17\%). There are no missing values in any column. Most transaction
amounts are small (under \$500 for both classes), but legitimate
transactions have outliers up to \$25,691 while fraud transactions never
exceed \$2,125 (Figure~\ref{fig:amount}). Fraudulent transactions are
more evenly spread across hours, with 2AM having the highest fraud rate
at 1.713\% (Figure~\ref{fig:hour}).

\subsection{Preprocessing}

The preprocessing pipeline (\texttt{src/transform.py}) performs three
steps: (1) engineer the \texttt{Hour} feature from the raw \texttt{Time}
column, (2) scale the \texttt{Amount} column using \texttt{StandardScaler}
and save the scaler artifact to \texttt{models/scaler.pkl} for consistent
inference, and (3) drop the original \texttt{Time} and \texttt{Amount}
columns. Median values of $V_1$--$V_{28}$ are saved to
\texttt{models/v\_medians.json} for use during single-transaction
prediction.

\subsection{Exploratory Data Analysis}

\begin{figure}[h]
\centering
\includegraphics[width=\columnwidth]{amount_boxplot.png}
\caption{Transaction amount by class. Legitimate transactions have
outliers up to \$25,691; fraud stays below \$2,125.}
\label{fig:amount}
\end{figure}

\begin{figure}[h]
\centering
\includegraphics[width=\columnwidth]{fraud_percentage_by_hour.png}
\caption{Fraud percentage by hour of day. 2AM has the highest fraud
rate at 1.713\%.}
\label{fig:hour}
\end{figure}

% ---------------------------------------------------------------
\section{Methodology}

\subsection{Evaluation Metrics}

Since the dataset is severely imbalanced, accuracy alone is misleading
--- a model predicting all transactions as legitimate achieves 99.83\%
accuracy. Instead, the following metrics are used:

\textbf{Precision} measures how many flagged transactions are actually
fraud:
\[
  \text{Precision} = \frac{TP}{TP + FP}
\]

\textbf{Recall} measures how many actual fraud cases the model catches:
\[
  \text{Recall} = \frac{TP}{TP + FN}
\]

\textbf{F1-Score} balances precision and recall:
\[
  F_1 = 2 \cdot \frac{\text{Precision} \cdot \text{Recall}}
                     {\text{Precision} + \text{Recall}}
\]

\textbf{AUC-ROC} measures the model's ability to distinguish between
classes across all thresholds (0.5 = random, 1.0 = perfect).

\subsection{Models}

\textbf{Baseline --- Logistic Regression:} A linear model that estimates
fraud probability using the sigmoid function:
\[
  P(y=1 \mid x) = \frac{1}{1 + e^{-(\beta_0 + \beta^T x)}}
\]

\textbf{LinearSVC:} A linear support vector classifier that finds the
maximum-margin hyperplane by minimizing:
\[
  \frac{1}{2}\|w\|^2 + C \sum_{i}
  \max\!\left(0,\ 1 - y_i(w^T x_i + b)\right)
\]

\textbf{Random Forest:} An ensemble of decision trees trained on
bootstrap samples. The final prediction is the majority vote across all
trees, which reduces overfitting and handles class imbalance.

\textbf{XGBoost:} A gradient boosting method that builds trees
sequentially, each correcting the residual errors of the previous trees
using a regularized objective function.

\textbf{LightGBM:} A gradient boosting method using leaf-wise tree
growth instead of level-wise, making it faster on large datasets.

\subsection{Training and Validation}

The data is split 80/20 into training (227,845 samples) and test
(56,962 samples). SMOTE is applied only to the training set, balancing
both classes to 227,451 samples each. The test set remains untouched to
simulate real-world conditions. All five models are trained on the
balanced training data and evaluated on the original imbalanced test set.

\begin{table}[h]
\centering
\caption{Model comparison on the test set.}
\label{tab:results}
\vskip 0.1in
\begin{small}
\begin{tabular}{lcccc}
\toprule
Model & Prec. & Recall & F1 & AUC \\
\midrule
Logistic Reg.  & 0.05 & 0.92 & 0.10 & 0.9715 \\
LinearSVC      & 0.05 & 0.92 & 0.10 & 0.9742 \\
Random Forest  & \textbf{0.86} & 0.83 & \textbf{0.84} & \textbf{0.9845} \\
XGBoost        & 0.54 & 0.87 & 0.66 & 0.9739 \\
LightGBM       & 0.10 & 0.91 & 0.19 & 0.9802 \\
\bottomrule
\end{tabular}
\end{small}
\end{table}

Random Forest achieves the best overall performance with the highest
precision (0.86), F1-score (0.84), and AUC-ROC (0.9845).

% ---------------------------------------------------------------
\bibliography{references}
\bibliographystyle{plainnat}

\end{document}